\documentclass[11pt]{article}
% ======================================================================
%  examstyle.tex — EPFL-style exam layout for the weekly Time Series exams
%  Shared by every Exam_Week_NN(.tex / _Solutions.tex). Compiles with pdflatex.
% ======================================================================
\usepackage[utf8]{inputenc}
\usepackage[T1]{fontenc}
\usepackage[english]{babel}
\usepackage[a4paper,margin=2.4cm]{geometry}
\usepackage{amsmath,amssymb,amsthm,mathtools}
\usepackage{bm}
\usepackage{enumitem}
\usepackage{array}

\setlength{\parindent}{0pt}
\setlength{\parskip}{0.5em}
\emergencystretch=3em
\allowdisplaybreaks[1]

% ---------- math macros (same names as the rest of the project) ----------
\newcommand{\E}{\mathbb{E}}
\DeclareMathOperator{\Var}{Var}
\DeclareMathOperator{\Cov}{Cov}
\DeclareMathOperator{\Corr}{Corr}
\DeclareMathOperator*{\argmin}{arg\,min}
\DeclareMathOperator{\MSE}{MSE}
\DeclareMathOperator{\trace}{tr}
\newcommand{\Reals}{\mathbb{R}}
\newcommand{\Ints}{\mathbb{Z}}
\newcommand{\Comp}{\mathbb{C}}
\newcommand{\Nat}{\mathbb{N}}
\newcommand{\Prob}{\mathbb{P}}
\newcommand{\Tr}{^{\mathsf{T}}}
\newcommand{\conj}[1]{\overline{#1}}
\newcommand{\abs}[1]{\left\lvert #1 \right\rvert}
\newcommand{\norm}[1]{\left\lVert #1 \right\rVert}
\newcommand{\ind}{\mathbf{1}}
\newcommand{\eps}{\varepsilon}
\newcommand{\diff}{\nabla}
\newcommand{\acvs}{\gamma}
\newcommand{\acf}{\rho}
\newcommand{\pacf}{\alpha}
\newcommand{\sdf}{\mathcal{S}}
\newcommand{\Bsh}{B}
\newcommand{\iid}{\overset{\text{iid}}{\sim}}

\setlist[enumerate]{leftmargin=2em,topsep=2pt,itemsep=4pt}
\setlist[itemize]{leftmargin=1.6em,topsep=2pt,itemsep=2pt}

\newcounter{exo}

% ---------- exam cover header :  \examheader{exam title}{duration} ----------
\newcommand{\examheader}[2]{%
  \thispagestyle{empty}
  \begin{center}
    {\Large\bfseries Time Series}\\[3pt]
    Spring Semester 2026, EPFL\\
    \'Ecole Polytechnique F\'ed\'erale de Lausanne\\
    Prof.\ Dr.\ Sofia C.\ Olhede\\[7pt]
    \hrule height 0.8pt
    \vspace{9pt}
    {\large\bfseries Practice Exam --- #1}\\[4pt]
    {\normalsize Duration: #2 \quad$\cdot$\quad No calculator \quad$\cdot$\quad Answer on paper}\\
    \vspace{8pt}
    \hrule height 0.8pt
  \end{center}
  \vspace{14pt}
  \begin{tabular}{@{}p{8.2cm}@{}p{8cm}@{}}
    Name:\enspace\hrulefill & Surname:\enspace\hrulefill
  \end{tabular}\\[14pt]
  SCIPER (Student Card Number):\enspace\hrulefill
  \vspace{16pt}
}

% ---------- points table (fixed: 4 problems) ----------
\newcommand{\pointstable}{%
  \begin{center}
  \renewcommand{\arraystretch}{1.5}
  \begin{tabular}{|p{5cm}|p{4cm}|}
    \hline
    \textbf{Exercise} & \textbf{Points}\\\hline
    1 & \\\hline
    2 & \\\hline
    3 & \\\hline
    4 & \\\hline
    \textbf{Total} & \\\hline
  \end{tabular}
  \end{center}
}

% ---------- problem heading (exam: one problem per page) ----------
\newcommand{\problem}[1]{%
  \clearpage
  \stepcounter{exo}%
  \begin{center}\textbf{\large Problem \theexo\ (#1 points)}\end{center}
  \vspace{4pt}
}
% ---------- answer space: fill the statement page + one blank answer page ----------
% Put \answerpage at the END of each problem so every exercise gets its own page(s).
\newcommand{\answerpage}{\par\vfill\newpage\null\newpage}

% ---------- solutions document ----------
\newcommand{\solheader}[1]{%
  \begin{center}
    {\Large\bfseries Time Series --- Solutions}\\[3pt]
    {\large Practice Exam --- #1}\\[2pt]
    EPFL, MATH-342 \quad$\cdot$\quad Prof.\ S.\ C.\ Olhede\\[5pt]
    \hrule height 0.8pt
  \end{center}
  \vspace{10pt}
}
\newcommand{\solproblem}[1]{%
  \bigskip
  \stepcounter{exo}%
  {\large\bfseries Problem \theexo\ (#1 points)}\par\nobreak\vspace{2pt}
}
% Rappel de l'énoncé avant la solution (dans les corrigés)
\newcommand{\statementstart}{\par\vspace{1pt}{\bfseries\itshape Statement.}\par\nobreak\begingroup\small\itshape}
\newcommand{\statementend}{\par\endgroup\vspace{5pt}\hrule\vspace{6pt}{\bfseries Solution.}\par\nobreak\vspace{2pt}}

% --- local macros not in examstyle ---
\DeclareMathOperator{\diag}{diag}
\begin{document}

\examheader{Week 10: Forecasting}{60 minutes}
\pointstable
\begin{center}\itshape Justify every answer. Throughout, $\eps_t$ is a mean-zero white noise of variance $\sigma^2$.\end{center}

% ---------------------------------------------------------------
\problem{5}
\textbf{(Forecasting a causal AR(1): point forecast, error variance, interval.)}
Let $\{X_t\}$ be the zero-mean Gaussian AR(1) process
\[
X_t=\phi X_{t-1}+\eps_t,\qquad \abs\phi<1,\quad \eps_t\iid\mathcal N(0,\sigma^2),
\]
and write $X^{\,n}_{n+h}$ for the best linear predictor of $X_{n+h}$ based on
$X_1,\dots,X_n$, with $h$-step forecast error $e_n(h)=X_{n+h}-X^{\,n}_{n+h}$.

\textbf{(a)} Give the $\mathrm{MA}(\infty)$ weights $\{\psi_j\}$ of the process. Using that the process is Gaussian (so the best linear predictor equals the conditional expectation), show that the point forecast is $X^{\,n}_{n+h}=\phi^{\,h}X_n$, and state its limit as $h\to\infty$. \hfill(1.5 pts)

\textbf{(b)} Show that the $h$-step forecast-error variance is
\[
\Var\bigl(e_n(h)\bigr)=\sigma^2\,\frac{1-\phi^{2h}}{1-\phi^2},
\]
and verify that $\Var(e_n(h))\to\acvs_0$ (the marginal variance) as $h\to\infty$. \hfill(2 pts)

\textbf{(c)} Now fix $\phi=0.6$ and $\sigma^2=4$. Report $X^{\,n}_{n+2}$ and the numerical forecast-error variances $\Var(e_n(1)),\Var(e_n(2))$, then give the $95\%$ Gaussian prediction interval for $X_{n+2}$ (use $z_{0.975}=1.96$). \hfill(1.5 pts)
\answerpage

% ---------------------------------------------------------------
\problem{5}
\textbf{(The prediction equations, and one-step AR(2) forecasts.)} Let $\{X_t\}$ be a zero-mean second-order stationary process with autocovariance $\{\acvs_\tau\}$, and let
\[
X^{\,n}_{n+h}=\sum_{j=1}^{n}\beta_j X_j
\]
be its best linear predictor of $X_{n+h}$ based on $X=(X_1,\dots,X_n)\Tr$.

\textbf{(a)} Starting from the prediction mean square error $S(\beta)=\E\bigl[(X_{n+h}-\beta\Tr X)^2\bigr]$, derive the \emph{prediction equations} $\Gamma_n\beta=\acvs_{[h]}$, where $\Gamma_n=[\acvs_{i-j}]_{i,j=1}^n$ and $\acvs_{[h]}=(\acvs_{n+h-1},\dots,\acvs_h)\Tr$, and show that this is equivalent to the orthogonality condition $\E[(X_{n+h}-X^{\,n}_{n+h})X_k]=0$ for $k=1,\dots,n$. \hfill(2 pts)

\textbf{(b)} Assuming $\Gamma_n$ invertible, deduce the closed forms
\[
X^{\,n}_{n+h}=\acvs_{[h]}\Tr\,\Gamma_n^{-1}X,
\qquad
P^{\,n}_{n+h}=\acvs_0-\acvs_{[h]}\Tr\,\Gamma_n^{-1}\acvs_{[h]} . \hfill(1.5\text{ pts})
\]

\textbf{(c)} Let $\{X_t\}$ be the causal AR(2) process $X_t=\phi_1 X_{t-1}+\phi_2 X_{t-2}+\eps_t$. Use the result of part (a) to show that the best one-step predictors based on $X_1$ and on $X_1,X_2$ are
\[
X^{\,1}_2=\frac{\acvs_1}{\acvs_0}X_1=\acf_1 X_1,
\qquad
X^{\,2}_3=\phi_1 X_2+\phi_2 X_1 .
\]
(For the second, you need \emph{not} invert $\Gamma_2$.) \hfill(1.5 pts)
\answerpage

% ---------------------------------------------------------------
\problem{5}
\textbf{(The causal-ARMA predictor via $\psi$-weights, and a worked ARMA(1,1).)}

\textbf{(a)} For a causal ARMA process with linear representation $X_t=\sum_{j=0}^\infty\psi_j\eps_{t-j}$, prove that the best linear predictor of $X_{n+h}$ based on $X_1,\dots,X_n$ is
\[
X^{\,n}_{n+h}=\sum_{j=h}^{\infty}\psi_j\,\eps_{n+h-j},
\qquad\text{with}\qquad
P^{\,n}_{n+h}=\sigma^2\sum_{j=0}^{h-1}\psi_j^2 .
\]
(\emph{Hint:} any linear predictor based on $X_1,\dots,X_n$ can be written $\sum_{j\ge0}c_j\eps_{n-j}$; expand the MSE and minimise over the $c_j$.) \hfill(2 pts)

\textbf{(b)} Consider the causal, invertible ARMA(1,1) process
\[
X_t=\tfrac12 X_{t-1}+\eps_t+\tfrac{3}{10}\eps_{t-1},\qquad \sigma^2=2 .
\]
Compute the first three $\psi$-weights $\psi_0,\psi_1,\psi_2$ (from $\Theta(B)/\Phi(B)$), and give the point forecasts $X^{\,n}_{n+1}$ and $X^{\,n}_{n+2}$ in terms of $X_n$ and the residual $\hat\eps_n$. \hfill(1.5 pts)

\textbf{(c)} Report the forecast-error variances $P^{\,n}_{n+1},P^{\,n}_{n+2},P^{\,n}_{n+3}$ as exact numbers, and explain in one sentence the value $\acvs_0$ they approach as $h\to\infty$ (you may use $\acvs_0=\sigma^2(1+2\phi\theta+\theta^2)/(1-\phi^2)$ for an ARMA(1,1)). \hfill(1.5 pts)
\answerpage

% ---------------------------------------------------------------
\problem{5}
\textbf{(Revision of Week 9: VAR stability and the diagonal VAR(1) covariance.)}

\textbf{(a)} Consider the bivariate $\mathrm{VAR}(2)$ process $X_t=\Phi_1X_{t-1}+\Phi_2X_{t-2}+\eps_t$ with
\[
\Phi_1=\begin{pmatrix}0.4 & 0\\[2pt] 0.5 & 0.9\end{pmatrix},\qquad
\Phi_2=\begin{pmatrix}0 & 0\\[2pt] 0.3 & -0.2\end{pmatrix}.
\]
Form the reverse characteristic polynomial $\det\bigl(I_2-\Phi_1 z-\Phi_2 z^2\bigr)$, exploit the lower-triangular structure to factor it, find all its roots, and decide whether the VAR(2) is stable. State the companion-matrix ($F$) form of the stability test and explain why it gives the same conclusion. \hfill(3 pts)

\textbf{(b)} Now take the \emph{diagonal} $\mathrm{VAR}(1)$ process $X_t=\Phi\,X_{t-1}+\eps_t$ with
\[
\Phi=\diag\!\bigl(\tfrac12,\tfrac14\bigr),\qquad
\{\eps_t\}\sim\mathrm{WN}(0,I_2).
\]
Write its $\mathrm{MA}(\infty)$ representation and compute the stationary covariance $\Gamma_0$ and the lag-$1$ matrix $\Gamma_1$ in closed form (exact fractions). \hfill(1.5 pts)

\textbf{(c)} For the process of part (b), what is the coherence $r_{1,2}(f)$ between the two channels at every frequency, and why? \hfill(0.5 pts)
\answerpage

\end{document}
